% PDF-derived LaTeX transcription. Original manuscript source was unavailable.
% Published PDF SHA-256: e0f88fa10165439149dd45d41b8adbbfe35afd7d90027137cb2c7f1ef0248e63
% Extracted text SHA-256: 638379946e4222d6dd7a9809db621c4164361622f8bfea61fb12417856a3e97e
\documentclass[11pt]{article}
\usepackage[margin=1in]{geometry}
\usepackage{amsmath,amssymb}
\setlength{\emergencystretch}{3em}
\title{Source-faithful static validation of archived CALVADOS-2 interactions: parameter equality, contact-range semantics, and backend-dependent test outcomes}
\author{}
\date{}
\begin{document}
\maketitle
\noindent\textbf{PDF-derived LaTeX transcription.} This editable text follows the PDF extraction. Refer to the published PDF for figures, table layouts, and mathematical notation.

% Extracted chunk 1

Source-faithful static validation of archived CAL V ADOS-2
interactions: parameter equality, contact-range semantics, and
backend-dependent test outcomes

\section*{Abstract}

Agreement between residue parameter tables does not establish equality of implemented
Hamiltonians. We examine a frozen comparison between an archived CAL V ADOS-2 direct-
coexistence source fixture and a saved OpenABC 1.0.9 comparison. Although the comparison
reports exact numerical parameter parity, an additional 4\ensuremath{\sigma} energy gate shortens four of 210
unordered canonical residue-pair interactions relative to the fixture’s fixed 2.0 nm contact cut-
off. A corrected contact-only implementation passes its original tolerances on saved CPU and
CUDA arrays comprising five selected pair types and 23 configurations. Independent local
reanalysis reproduces six error statistics, array and contract identities, residue-table mixing,
and an analytic energy/derivative reference without rerunning OpenMM. A separate, broader
static harness covers 1,470 contact configurations, five bond configurations, and 30 electrostatic
configurations. Its saved Reference report passes all sections, whereas its CUDA report fails
contact-force and bond tolerances while passing electrostatic/terminal checks. Independent
report-consistency checks preserve, rather than resolve, these failures. The results distinguish
parameter agreement, functional equivalence, fixture-specific numerical agreement, and physical
validation; integration, ensembles, transferability, and GPU performance remain unestablished.

\section*{1 Introduction and contributions}

Residue-level simulations require a precisely specified interaction model before numerical agreement
can be interpreted physically. A residue table supplies coefficients, but the evaluated potential also
depends on parameter mixing, branch definitions, cutoff placement, energy shifts, exclusions, and
unit conventions. Backend selection and force evaluation introduce another level of comparison.
Consequently, neither matching coefficients nor agreement on a small collection of configurations
establishes equality of complete Hamiltonians.
This study addresses that distinction through a bounded historical CAL V ADOS-2 compari-
son. The reference fixture is a captured direct-coexistence script associated with Zenodo archive
10.5281/zenodo.13754000, version 2.0. The candidate comparison identifies OpenABC 1.0.9. These
identities delimit the inference: the findings concern the saved comparison, not current OpenABC
releases or every CAL V ADOS implementation.
The focused contribution is threefold. First, we explain how identical reported residue parame-
ters coexist with different contact ranges. Second, we quantify what is reproduced from a corrected
23-configuration contact fixture. Third, we compare those results with the contrasting outcomes of
a separate, more comprehensive static harness, including its CUDA failures. This bounded imple-
mentation comparison establishes the scope of the available measurements; it does not introduce a
new force field.
1


\section*{2 Related work}

Tesei and colleagues demonstrated that interaction range affects CAL V ADOS phase-separation
predictions more strongly than single-chain compaction. Their model-development analysis distin-
guishes optimization at 2.4 nm from simulations evaluated at 2.0 nm [ 1]. This context motivates
checking range semantics; it does not independently authenticate the private comparison examined
here. In particular, the fixed 2.0 nm cutoff below belongs to the captured source fixture, not a
universal definition of CAL V ADOS.
OpenABC provides Python-based access to multiple coarse-grained force fields through OpenMM.
Its publication includes cross-engine energy comparisons and separate ensemble and performance
assessments [ 2]. Those results establish relevant implementation context, not the outcome of this
archived CAL V ADOS-2/OpenABC 1.0.9 comparison. We therefore distinguish published software
capabilities from the source-specific evidence and numerical contracts analyzed here.

\section*{3 Scientific methods and validation protocols}

Evidence and retrospective analysis
This is a retrospective analysis of captured code, reports, contracts, configuration lists, and small
saved arrays. No additional molecular-dynamics simulations or OpenMM backend executions were
performed in the local reanalysis. Supplied source code was inspected, not executed. Two inde-
pendent local analyses examined the saved numerical evidence: one recalculated the small-array
results; the other performed 37 checks of report arithmetic, aggregation of pairwise maxima, thresh-
old decisions, and links to declared captured sources. Neither analysis constitutes a fresh simulator
reproduction.
The archived direct-coexistence source explicitly sets the contact cutoff and shift distance to 2.0
nm and the electrostatic cutoff and shift distance to 4.0 nm. It excludes directly bonded neighbors
from both nonbonded terms and enables periodic boundaries. The static pair harnesses instead
use their default nonperiodic configurations. Thus, even broad pair coverage is not a periodic
multichain-system test. Static archive-member comparisons establish that the captured OpenABC
interaction source matches the corresponding member of the official OpenABC 1.0.9 package and
that the captured direct-coexistence script and residue table match members of Zenodo archive
10.5281/zenodo.13754000, version 2.0. The source identities and retrieval records are preserved in
Supplementary record S1. These byte-identity checks do not reproduce the saved backend execu-
tions.
Contact potential, mixing, and force convention
For residue types i, j, the captured preparation uses arithmetic mixing,
\ensuremath{\sigma}ij = \ensuremath{\sigma}i + \ensuremath{\sigma}j
2 , \ensuremath{\lambda} ij = \ensuremath{\lambda}i + \ensuremath{\lambda}j
2 , \ensuremath{\varepsilon} = 0.2 \ensuremath{\times} 4.184 = 0 .8368 kJ mol \ensuremath{-}1. (1)
2

With rm = 2 1/6\ensuremath{\sigma}ij, rc = 2 .0 nm, and uLJ(r) = 4 \ensuremath{\varepsilon}[(\ensuremath{\sigma}ij/r)12 \ensuremath{-} (\ensuremath{\sigma}ij/r)6], the source-faithful shifted
Ashbaugh–Hatch contact potential is
UAH(r) =
[U+F8F1]
[U+F8F4][U+F8F2]
[U+F8F4][U+F8F3]
uLJ(r) + (1 \ensuremath{-} \ensuremath{\lambda}ij)\ensuremath{\varepsilon} \ensuremath{-} \ensuremath{\lambda}ijuLJ(rc), 0 < r < r m,
\ensuremath{\lambda}ij[uLJ(r) \ensuremath{-} uLJ(rc)], rm \ensuremath{\leq} r < r c,
0, r \ensuremath{\geq} rc.
(2)
The identical cutoff-energy shift occurs in both branches. The repulsive-branch derivative is not
multiplied by hydropathy. Away from branch and cutoff boundaries,
U \ensuremath{\prime}
AH(r) = a(r)4\ensuremath{\varepsilon}
[
\ensuremath{-}
12\ensuremath{\sigma}12
ij
r13 +
6\ensuremath{\sigma}6
ij
r7
]
, a (r) =
\{
1, r < r m,
\ensuremath{\lambda}ij, r m < r < r c. (3)
For particles at (0, 0, 0) and (r, 0, 0), F1x = + U \ensuremath{\prime}(r) and F2x = \ensuremath{-}U \ensuremath{\prime}(r); transverse components
vanish. The shift makes the energy continuous at the cutoff, but is not a force-shifting prescription.
The tests do not resolve the force exactly at rc.
The saved OpenABC comparison reports the same 2.0 nm neighbor-list limit and shift distance,
plus an additional multiplicative step(4\ensuremath{\sigma}ij \ensuremath{-} r) energy gate. Hence its effective range is the smaller
of these two distances. This is a functional difference even when all table entries agree.
Small corrected contact fixture
The corrected contact implementation removes that additional gate. Its saved fixture selects A–G,
C–G, G–G, G–S, and Y–W: five pair types, not all 210. The first four receive five distance entries
around their gates and near 1.999 nm; Y–W receives 1.70, 1.90, and 1.999 nm. The denominator
is 23 configuration entries, including a repeated C–G entry at 1.999 nm. All are on the attractive
branch. The fixture does not cover repulsion, exclusions, bonds, or electrostatics.
The original harness requests a single CPU thread or CUDA double precision and evaluates static
states without advancing its Verlet integrator. Saved energies have shape [23] and forces [23, 2, 3].
Independent local decoding used pickle-disabled array loading, checked exact array identities against
the contracts, verified the two captured residue tables and arithmetic mixing, and recalculated
maximum absolute, RMS, and relative errors for energy and force. Energy relative error uses a
pointwise denominator max(|Eref |, 10\ensuremath{-}12 kJ mol\ensuremath{-}1); force relative error is a global flattened L2
ratio. Force RMS includes all Cartesian components. Energy RMS is descriptive, not a frozen
acceptance gate.
Separate broader static harness
The subsequent full-semantics harness uses a different captured implementation source and a dif-
ferent test program. Each of 210 unordered residue pairs is evaluated at 0.90rm, rm \ensuremath{-} 10\ensuremath{-}4 nm,
rm + 10 \ensuremath{-}4 nm, 1.0, 1.5, 1.999, and 2.001 nm, yielding 1,470 contact configurations. A separate
excluded contact pair is evaluated at 0.7 nm. Five bond distances are 0.20, 0.3799, 0.38, 0.3801,
and 0.55 nm.
The captured harmonic convention is
Ub(r) = 1
2 kb(r \ensuremath{-} r0)2, r 0 = 0.38 nm , k b = 8033 kJ mol \ensuremath{-}1 nm\ensuremath{-}2. (4)
3

Sequence preparation assigns histidine charge [1 + 10 pH\ensuremath{-}6]\ensuremath{-}1, adds +1/ \ensuremath{-} 1 elementary charge to
N/C termini, and adds +2/ + 16 amu to their masses. The helper rejects one-residue chains. Its
shifted screened electrostatic potential is
UDH(r) = qiqj\ensuremath{\ell}BkBT
[e\ensuremath{-}\ensuremath{\kappa}r
r \ensuremath{-} e\ensuremath{-}\ensuremath{\kappa}re
re
]
, r < r e = 4.0 nm , (5)
with zero interaction beyond the cutoff. Here charges are elementary-charge numbers; the captured
temperature-dependent dielectric calculation determines \ensuremath{\ell}B, and ionic strength determines \ensuremath{\kappa}.
Five sequence/pH/temperature/salt cases are HGH/5.0/280 K/0.05 M, RHK/6.0/298 K/0.15
M, DHE/7.0/310 K/0.50 M, AHHHA/7.4/323 K/0.15 M, and KGE/9.0/350 K/0.05 M. Terminal
pairs, plus one internal pair for AHHHA, are evaluated at 0.5, 1.0, 2.0, 3.999, and 4.001 nm: 30
electrostatic configurations. The RHK and DHE terminal pairs have zero charge products and
therefore test zero-interaction behavior. These are pair evaluations, not sequence trajectories. The
harness requests CUDA double precision; Reference uses its default properties. Original thresholds
are retained throughout.

\section*{4 Experimental results}

Parameter parity does not repair the range mismatch
The saved comparison reports exact numeric parameter-table parity but four prematurely truncated
pairs among 210. Table 1 gives every affected pair. The additional gate removes up to 0.200 nm
of the intended range for G–G. It also terminates an interaction shifted at a different distance;
matching the nominal neighbor cutoff therefore does not establish matching evaluated energies or
forces.
Table 1: All affected pairs in the frozen comparison. Distances are in nm; the source cutoff is 2.000
nm for every row. Truncations are displayed at the precision of the gate distances.
Pair \ensuremath{\sigma}ij Additional gate 4\ensuremath{\sigma}ij Lost range
A–G 0.477 1.908 0.092
C–G 0.499 1.996 0.004
G–G 0.450 1.800 0.200
G–S 0.484 1.936 0.064
Reproduced small-array agreement
Both saved corrected contact datasets pass their original five gates (Table 2). Local recalculation
reproduces all six statistics to floating-point rounding. The independently formulated reference
matches saved reference energies exactly on these entries; its maximum force difference is 2.082 \ensuremath{\times}
10\ensuremath{-}17 kJ mol\ensuremath{-}1 nm\ensuremath{-}1. This verifies the saved reference and analytic derivative on the supplied
attractive-branch configurations, not a new backend execution.
The CPU maximum relative energy error is sensitive to small absolute discrepancies because the
shifted reference energy approaches zero near 2.0 nm. Absolute and relative gates assess different
aspects of agreement; no tolerance was changed during reanalysis.
4

Table 2: Small contact-only fixture: 23 entries on each platform. Energy errors are in kJ mol\ensuremath{-}1;
force errors in kJ mol\ensuremath{-}1 nm\ensuremath{-}1; relative errors are dimensionless. Values are the saved statistics
independently recalculated locally, displayed with up to four significant figures.
Statistic CPU CUDA Frozen limit
Energy maximum absolute 2.011 \ensuremath{\times} 10\ensuremath{-}9 1.232 \ensuremath{\times} 10\ensuremath{-}10 10\ensuremath{-}6
Energy RMS 6.751 \ensuremath{\times} 10\ensuremath{-}10 2.994 \ensuremath{\times} 10\ensuremath{-}11 Not gated
Energy maximum relative 2.295 \ensuremath{\times} 10\ensuremath{-}4 1.847 \ensuremath{\times} 10\ensuremath{-}7 5 \ensuremath{\times} 10\ensuremath{-}4
Force maximum absolute 1.083 \ensuremath{\times} 10\ensuremath{-}8 5.883 \ensuremath{\times} 10\ensuremath{-}10 10\ensuremath{-}5
Force RMS 1.83 \ensuremath{\times} 10\ensuremath{-}9 1.243 \ensuremath{\times} 10\ensuremath{-}10 10\ensuremath{-}6
Force relative L2 3.26 \ensuremath{\times} 10\ensuremath{-}7 2.214 \ensuremath{\times} 10\ensuremath{-}8 10\ensuremath{-}5
Broader Reference pass and CUDA failure
The broader saved Reference report passes contact, bond, and electrostatic/terminal sections. The
CUDA report fails overall: contact maximum and RMS force errors exceed their frozen limits, and
all three bond error statistics exceed theirs. Electrostatic/terminal checks pass. Table 3 retains
these outcomes and all reported aggregate residuals needed to interpret them.
Table 3: Separate full-semantics harness. Emax is in kJ mol\ensuremath{-}1; Fmax and FRMS are in kJ mol\ensuremath{-}1 nm\ensuremath{-}1.
Counts exclude the separate contact-exclusion probe. Error values are displayed with up to four
significant figures.
Section/platform Count Emax Fmax FRMS
Contact/Reference: PASS 1470 4.441 \ensuremath{\times} 10\ensuremath{-}16 5.329 \ensuremath{\times} 10\ensuremath{-}14 4.696 \ensuremath{\times} 10\ensuremath{-}15
Contact/CUDA: F AIL 1470 8.938 \ensuremath{\times} 10\ensuremath{-}7 3.16 \ensuremath{\times} 10\ensuremath{-}5 3.575 \ensuremath{\times} 10\ensuremath{-}6
Contact limits 10\ensuremath{-}6 10\ensuremath{-}5 10\ensuremath{-}6
Bond/Reference: PASS 5 0 0 0
Bond/CUDA: F AIL 5 6.895 \ensuremath{\times} 10\ensuremath{-}6 3.83 \ensuremath{\times} 10\ensuremath{-}5 2.211 \ensuremath{\times} 10\ensuremath{-}5
Bond limits 10\ensuremath{-}9 10\ensuremath{-}8 10\ensuremath{-}9
Electrostatic/Reference: PASS 30 1.776 \ensuremath{\times} 10\ensuremath{-}15 1.421 \ensuremath{\times} 10\ensuremath{-}14 1.504 \ensuremath{\times} 10\ensuremath{-}15
Electrostatic/CUDA: PASS 30 8.55 \ensuremath{\times} 10\ensuremath{-}7 2.385 \ensuremath{\times} 10\ensuremath{-}6 2.634 \ensuremath{\times} 10\ensuremath{-}7
Electrostatic limits 10\ensuremath{-}6 10\ensuremath{-}5 10\ensuremath{-}6
Contact relative force L2 errors are 7.615 \ensuremath{\times} 10\ensuremath{-}16 for Reference and 5.798 \ensuremath{\times} 10\ensuremath{-}7 for CUDA,
both below 10\ensuremath{-}5. Both reports record zero excluded-pair contact energy and force, against 10\ensuremath{-}15
limits in their respective units. Charge, mass, and electrostatic-constant discrepancies are also zero,
against declared 10\ensuremath{-}15 limits. The constants check is a numerical comparison of the source-defined
tuple, not a common-unit physical observable.
All 37 report-consistency checks pass. They confirm that section decisions follow the declared
thresholds, that the 210 pair maxima aggregate to the reported contact maxima, and that the
reports join to the declared captured sources. The CUDA contact maxima occur in the E–G
entry. These checks establish internal consistency of a failing report, not successful backend parity.
Missing full-semantics arrays prevent independent reconstruction of individual residuals or their
RMS values.
5


\section*{5 Discussion and limitations}

The comparison separates two different reasons that a favorable preliminary result may not support
a broader conclusion. Parameter equality fails to establish functional equality because the extra
distance gate changes the evaluated interaction. Subsequently, a corrected implementation can
agree on selected attractive-tail configurations while another implementation source and broader
harness fail their own numerical contract. Neither inference requires extrapolating static errors into
an unmeasured thermodynamic effect.
The small CUDA pass does not override the broader CUDA failure. The harnesses differ in
implementation source, tested distances, force terms, and acceptance contracts, particularly the
much tighter bond tolerances. The broader contact test includes repulsive configurations and
all canonical pairs; the narrow fixture does not. Conversely, the small fixture has array-level
reproduction that the broader reports lack. These complementary evidence levels should not be
collapsed into a single backend-validation label.
The cause of the broader CUDA residuals is unresolved. Although the captured harness requests
double precision, neither report arithmetic nor the supplied source establishes a causal explanation
for the discrepancies. Assigning them to precision would exceed the evidence. Likewise, passing
Reference tests does not prove universal exactness: zero residuals on five bond configurations are
finite-test observations.
Static energy and force evaluation tests neither integration stability nor equilibrium sampling.
No independent trajectories were analyzed in this study. Periodic multichain exclusions, accumu-
lated forces in dense systems, thermostat behavior, ensemble observables, phase coexistence, and
transferability require additional evidence. GPU throughput was not measured. The published
range-sensitivity results motivate attention to cutoffs but do not quantify the physical consequences
of these four pair-specific truncations.
The authenticated OpenABC 1.0.9 source confirms the additional gate in the cutoff-parameterized
Ashbaugh–Hatch function used in the frozen comparison. This supports the version-bound func-
tional discrepancy, not a claim about every configuration in that release. Current-release claims
would require a new, version-bound source comparison.

\section*{6 Conclusion}

For the archived fixed-2.0 nm CAL V ADOS-2 fixture, the saved OpenABC 1.0.9 comparison reports
matching parameters but four shortened contact ranges. Corrected small contact arrays repro-
duce their frozen acceptance results under independent local arithmetic and analytic checks. A
different broader static harness passes on Reference but fails contact-force and bond tolerances on
CUDA. These failures remain scientifically material even though report-consistency checks succeed.
Source-faithful validation therefore requires explicit separation of model semantics, fixture coverage,
numerical thresholds, backend evidence, and subsequent integration and ensemble validation.

\section*{7 Data and Code A vailability}

The supplied private evidence contains captured contact and full-semantics implementation code,
both harnesses, a parity evaluator, configuration lists, contracts, manifests, comparison reports,
residue-table artifacts, and four small contact NPZ arrays. The two independent local-check out-
6

puts are also supplied. Their inventory and limitations are identified in Supplementary record
S1. Access in this study was through the private supplied export; no public archive, reviewer-
access mechanism, or redistribution permission has been established. The export also contains
the authenticated OpenABC 1.0.9 interaction-source member and the primary archive identity and
retrieval records. The upstream package and CAL V ADOS archive are available at https://pypi.
org/project/openabc/1.0.9/ and https://zenodo.org/records/13754000, respectively; these
upstream resources do not provide access to the private numerical comparison. Raw full-semantics
energy/force arrays and usable trajectory datasets are absent. Upstream publications identify their
own code resources, but those resources do not constitute public access to this private compari-
son. Further sharing of captured code, tables, arrays, and recipes is subject to source-access and
applicable redistribution/licensing restrictions.

\section*{References}

[1] G. Tesei and K. Lindorff-Larsen. Improved predictions of phase behaviour of intrinsically dis-
ordered proteins by tuning the interaction range. Open Research Europe 2:94, version 2, 2023
(2022 collection). doi:10.12688/openreseurope.14967.2.
[2] S. Liu, C. Wang, A. P. Latham, X. Ding, and B. Zhang. OpenABC enables flexible, simplified,
and efficient GPU accelerated simulations of biomolecular condensates. PLoS Computational
Biology 19(9):e1011442, 2023. doi:10.1371/journal.pcbi.1011442.
7

\end{document}
